% Ad Familiares 13.47 --- headnote
% ad-familiares-13-47, ca. end of March 51 BC, Rome

\textit{Cicero to L. Silius, written at Rome around the end of
March 51 BC (Perseus: \textit{ut videtur, ex. m. Martio} ---
the year noted as 701 in the manuscript is a slip for 703).
A brief recommendation of Egnatius of Sidicinum: the kind of
formal letter Cicero produced in large numbers, especially as
he prepared to leave Rome for his Cilician command. Egnatius
will appear again in Cicero's correspondence as a man of
business closely tied to him; here he is being commended to
Silius's good offices.}

\textit{The shape is the familiar one. There is no need to
recommend a man Silius already loves; but Cicero will write all
the same, so that Silius knows Egnatius is loved, not merely
esteemed. Then the request, and then the unexplained domestic
shadow that lifts the letter out of the formula: \textit{illa
nostra ceciderunt} --- ``those things of ours have fallen
through.'' Some shared concern, or some business, has come to
nothing; the particulars are between the two men. Cicero
reaches for the common consolation, ``what if this is for the
better,'' and turns at once to the closing. \textit{Sed haec
coram} --- ``but these things in person'' --- is the
correspondence's stock formula for leaving a delicate subject
unwritten when the writer expects to see the addressee soon.
The whole sits within a few lines and yet, like many of the
\textit{Familiares} 13 letters of recommendation, opens a
window onto a private affair the modern reader can only watch
slip closed.}
